\chapter{Análisis del problema}\label{cap:analisis}

\section{Descripción del problema}

El problema que aborda GlucoMate nace de una necesidad doble. Por un lado, las personas
con diabetes tipo 1 necesitan consultar su glucosa y contextualizarla continuamente con
otros factores clínicos como la insulina administrada, las comidas y la actividad
deportiva. Por otro, las soluciones disponibles suelen fragmentar esa información,
limitar la personalización del tratamiento o no ofrecer una capa de apoyo pedagógico que
ayude a interpretar los datos.

Desde el punto de vista técnico, el proyecto debe integrar en una única solución móvil la
captura de datos del sensor, su persistencia, su visualización histórica y la generación
de recomendaciones adaptadas al perfil del usuario.

El análisis del proyecto real confirma además varias restricciones de partida. La primera
es la dependencia del ecosistema Android para el acceso a NFC y para la integración con la
cadena de herramientas elegida. La segunda es la necesidad de apoyarse en una aplicación
intermedia, Juggluco, para disponer del dato glucémico descifrado. La tercera es la
separación entre frontend y backend, necesaria para no exponer en el cliente las
credenciales de base de datos ni la clave del servicio de IA.

\section{Requisitos funcionales}

Entre los requisitos funcionales principales del sistema se encuentran los siguientes:

\begin{itemize}
    \item[\textbf{RF-001:}] Leer lecturas de glucosa del sensor Freestyle Libre 2 Plus
          desde un dispositivo Android.
    \item[\textbf{RF-002:}] Solicitar el dato glucémico reciente a un servicio puente
          accesible mediante HTTP local.
    \item[\textbf{RF-003:}] Mostrar el valor actual de glucosa y su tendencia de forma
          clara e inmediata.
    \item[\textbf{RF-004:}] Representar la evolución histórica reciente mediante una
          gráfica temporal.
    \item[\textbf{RF-005:}] Registrar eventos clínicos relevantes: comidas, insulina y
          actividad física.
    \item[\textbf{RF-006:}] Gestionar un perfil configurable con parámetros terapéuticos
          del usuario.
    \item[\textbf{RF-007:}] Persistir los datos en una base de datos para su posterior
          consulta.
    \item[\textbf{RF-008:}] Generar recomendaciones orientativas a través de un agente
          de IA.
    \item[\textbf{RF-009:}] Recuperar catálogos auxiliares, como insulinas y deportes,
          para facilitar la entrada de datos.
    \item[\textbf{RF-010:}] Mantener sincronizados frontend, backend y base de datos tras
          cada operación de lectura o registro.
\end{itemize}

\section{Requisitos no funcionales}

Los requisitos no funcionales más relevantes son:

\begin{itemize}
    \item[\textbf{RNF-01:}] La aplicación debe poder ser utilizada por personas no
          expertas, con una interfaz directa y comprensible.
    \item[\textbf{RNF-02:}] La lectura, almacenamiento y renderizado de datos deben
          producirse sin demoras perceptibles para el usuario.
    \item[\textbf{RNF-03:}] El sistema debe gestionar errores de red, NFC o backend sin
          comprometer la estabilidad general de la aplicación.
    \item[\textbf{RNF-04:}] La solución debe organizarse en módulos separados para
          facilitar su evolución futura.
    \item[\textbf{RNF-05:}] Las credenciales y claves de servicios externos no deben
          distribuirse dentro del cliente móvil.
\end{itemize}

\section{Requisitos de información}

Para operar correctamente, GlucoMate necesita manejar al menos los siguientes bloques de
información:

\begin{itemize}
    \item[\textbf{RI-01:}] Lecturas de glucosa con su correspondiente marca temporal.
    \item[\textbf{RI-02:}] Datos del perfil terapéutico del usuario.
    \item[\textbf{RI-03:}] Registros de administración de insulina.
    \item[\textbf{RI-04:}] Registros de comidas y su información asociada.
    \item[\textbf{RI-05:}] Registros de actividad física y ejercicio.
    \item[\textbf{RI-06:}] Respuestas generadas por el agente de apoyo.
    \item[\textbf{RI-07:}] Formato homogéneo de los datos para poder cruzarlos en la
          visualización y en la lógica de recomendación.
\end{itemize}

En términos más concretos, el backend del proyecto trabaja con entidades equivalentes a
\texttt{glucose\_readings}, \texttt{user\_profile}, \texttt{insulins},
\texttt{insulin\_events}, \texttt{food\_events}, \texttt{sports} y
\texttt{exercise\_events}. Esta estructura confirma que la aplicación no se limita a una
mera lectura puntual del sensor, sino que construye una historia clínica reducida pero
operativa para personalizar la ayuda mostrada al usuario.

\section{Casos de uso}

Los casos de uso principales del sistema son: consultar la glucosa actual, revisar el
histórico, registrar una comida, registrar una dosis de insulina, registrar actividad
física, editar el perfil terapéutico y solicitar una recomendación al agente de IA.

El actor principal es la persona con diabetes que utiliza la aplicación desde su teléfono
móvil. Como actores secundarios pueden considerarse el backend de GlucoMate, la base de
datos PostgreSQL, Juggluco como servicio puente y el proveedor externo del modelo de IA.
La mayor parte de los casos de uso combinan interacción local en el móvil con llamadas de
red a servicios de apoyo.

\subsection{Consultar la glucosa actual}

\textbf{Actor:} usuario con diabetes.

\textbf{Flujo principal:} el usuario accede a la pantalla principal de la aplicación y
solicita la lectura más reciente del sensor. La aplicación obtiene el dato glucémico a
través del servicio puente, lo procesa y muestra en pantalla el valor actual junto con su
tendencia. Este caso de uso concentra la interacción esencial del sistema, ya que permite
al usuario conocer de forma inmediata su situación glucémica actual.

\subsection{Revisar el histórico}

\textbf{Actor:} usuario con diabetes.

\textbf{Flujo principal:} el usuario abre la vista histórica desde la aplicación. El
sistema recupera las lecturas almacenadas en el backend y las representa en una gráfica
temporal junto con los eventos clínicos asociados. De este modo, el usuario puede
interpretar la evolución reciente de la glucosa en relación con comidas, insulina y
actividad física.

\subsection{Registrar una comida}

\textbf{Actor:} usuario con diabetes.

\textbf{Flujo principal:} el usuario selecciona la opción de registrar una ingesta,
introduce la información relevante y confirma la operación. La aplicación envía el evento
al backend, que lo persiste en la base de datos y lo deja disponible para su posterior
visualización. El registro queda integrado con el resto de la información clínica del
usuario.

\subsection{Registrar una dosis de insulina}

\textbf{Actor:} usuario con diabetes.

\textbf{Flujo principal:} el usuario accede al formulario de insulina, selecciona el tipo
de insulina, indica la dosis administrada y guarda el registro. La aplicación comunica la
información al backend, que la almacena de forma estructurada para su uso posterior en la
visualización y en la lógica de recomendación. Así, el sistema puede contextualizar la
curva glucémica con la pauta de insulina aplicada.

\subsection{Registrar actividad física}

\textbf{Actor:} usuario con diabetes.

\textbf{Flujo principal:} el usuario abre la opción de actividad física, selecciona el
deporte o tipo de ejercicio realizado, completa los datos requeridos y confirma el
registro. El sistema guarda la información en el backend y la incorpora a la historia
temporal del usuario. Esto permite analizar la posible influencia del ejercicio sobre la
evolución glucémica posterior.

\subsection{Editar el perfil terapéutico}

\textbf{Actor:} usuario con diabetes.

\textbf{Flujo principal:} el usuario entra en la pantalla de perfil y actualiza sus
parámetros terapéuticos, como objetivos glucémicos, sensibilidad, ratio de hidratos o
configuración de insulinas. Tras guardar los cambios, la aplicación envía la información
al backend, donde queda persistida como referencia clínica principal del usuario. Este
perfil sirve como base para personalizar la visualización y las recomendaciones.

\subsection{Solicitar una recomendación al agente de IA}

\textbf{Actor:} usuario con diabetes.

\textbf{Flujo principal:} el usuario solicita ayuda desde la interfaz de la aplicación.
El sistema recopila el contexto clínico disponible, incluyendo glucosa reciente, perfil
terapéutico y eventos registrados, y lo envía al backend para construir la consulta al
agente de IA. La respuesta generada se devuelve a la aplicación y se muestra al usuario en
forma de orientación interpretativa sobre su situación actual.

\section{Modelo de dominio}

El modelo de dominio de GlucoMate gira en torno a las entidades \textit{Lectura de
Glucosa}, \textit{Usuario}, \textit{Perfil Terapéutico}, \textit{Evento de Insulina},
\textit{Comida}, \textit{Actividad Física} y \textit{Recomendación}. La relación entre
ellas se desarrolla con más detalle en el capítulo de diseño, donde se concreta también
su traducción al modelo relacional de la base de datos.

Una característica relevante del dominio es su fuerte dependencia temporal. Casi todas las
entidades poseen marcas de tiempo y se interpretan mejor como eventos dentro de una ventana
reciente: la glucosa de las últimas 12 horas, la insulina activa de las últimas 4 horas o
el ejercicio practicado dentro de la ventana de riesgo definida para cada deporte.